\subsection{Corte de vara}

El problema de corte de vara consiste en determinar cuál es la mejor forma de cortar una vara de longitud $n$, si cada posible longitud entera se puede vender a un precio específico. El objetivo es maximizar los ingresos.

Las entradas del problema son entonces una longitud entera $n$ y una función $p$ que determina el precio que tendría cada pedazo. Para facilidad computacional, usualmente el problema se establece como:
\begin{itemize}
  \item Entradas: Un entero $n$ que representa la longitud de la vara y un arreglo $p$ de tamaño $n+1$, donde $p[i]$ es el precio de una vara de longitud $i$.
  \item Salida: Los ingresos máximos que se pueden obtener al cortar la vara.
\end{itemize}
En ese modelo, $p[0] = 0$ porque longitud $0$ representa la ausencia de vara (y por ende de venta).

Considérese la instancia del problema con $n=4$ y $p = [0, 1, 5, 8, 9]$. En ese caso:
\begin{itemize}
  \item Si no se corta la vara, se vende a $p[4]$, o sea a $\$9$.
  \item Si se corta en un pedazo de longitud $3$ y otro de longitud $1$, se venden a $\$8$ y $\$1$, respectivamente, para obtener ingresos de $\$9$.
  \item Si se corta en dos pedazos de longitud $2$, se venden a $\$10$.
  \item Si se corta en pedazos de longitud $2$, $1$ y $1$, se venden a $\$7$.
  \item Si se corta en cuatro pedazos de longitud $1$, se venden a $\$4$.
\end{itemize}
Entonces, la mejor forma de cortar la vara es en dos pedazos de longitud $2$, para obtener ingresos de $\$10$. Como $n$ es pequeño, esta instancia se resuelve por fuerza bruta, mirando todas las formas de cortar la vara y calculando los ingresos para cada una.

Para un valor arbitrario de $n$, hay muchas formas de cortar la vara. En cada unidad de longitud, hay dos alternativas: se corta o no. Entonces hay $2$ posibilidades para cada una de las $n-1$ unidades de longitud ($n-1$ porque en la longitud $n$ ya no se puede cortar), lo que da $\underbrace{2\cdot2\cdots2}_{n-1 \text{ veces}} = 2^{n-1}$ combinaciones totales.

El conteo anterior no contempla que algunas combinaciones están repetidas, porque el orden de los pedazos no importa: es lo mismo cortar en la primera unidad de longitud y obtener un pedazo de longitud $1$ y otro de longitud $n-1$, que en la penúltima unidad de longitud, obteniendo un pedazo de longitud $n-1$ y otro de longitud $1$. Esta observación tiene dos implicaciones:
\begin{itemize}
  \item Primero, reduce el número de combinaciones de $2^{n-1}$ a un valor dado por la función de partición de enteros, $\mathrm{e}^\mathrm{\pi}\sqrt{2n/3}/(4n\sqrt{3})$. Esta es una gran mejora, pero la cantidad de combinaciones sigue lejos de ser polinomial. (No se muestra la deducción de la función de partición porque es un problema de conteo mucho más difícil).
  \item Segundo, implica que no se debe pensar en cortar o no en cada unidad, sino que se puede pensar en cortar de forma que los pedazos queden en tamaño no decreciente. Eso igual va a contemplar todas las combinaciones: no contempla ($n-2$, $2$), ni ($n-2$, $1$, $1$); pero sí ($2$, $n-2$), también ($1$, $1$, $n-2$), etc., y como dan los mismos ingresos, con eso basta.
\end{itemize}

\subsubsection{Solución recursiva}

Para diseñar una solución recursiva, lo primero que hay que pensar es cómo dividir este problema en subproblemas. En este caso es relativamente fácil, porque dada una vara de longitud $n$, se querría evaluar:
\begin{itemize}
  \item El primer caso de cortar en la primera unidad de longitud. Eso arroja un pedazo de longitud $1$ y otro de longitud $n-1$. Ese último pedazo de longitud $n-1$ se puede seguir cortando (si $n$ es grande) y entonces se tiene el mismo problema: cuál es la forma de cortarlo para maximizar los ingresos.
  \item El segundo caso es cortar en la segunda unidad de longitud y es análogo: el pedazo de longitud $2$ se vende a $p[2]$, pero falta evaluar cómo cortar el pedazo de longitud $n-2$ de forma óptima.
\end{itemize}
Con eso es fácil construir una implementación recursiva, \inline{corte_de_vara_recursivo} que evalúe el primer caso como \inline{p[1] + corte_de_vara_recursivo(n-1, p)}; el segundo, como \inline{p[2] + corte_de_vara_recursivo(n-2, p)} y así hasta \inline{p[n] + corte_de_vara_recursivo(0, p)}, donde ese último llamado recursivo evidentemente es un caso base.
\begin{pseudocode}
función corte_de_vara_recursivo(n: entero, p: arreglo de enteros) entero
    if n <= 0 return p[0]

    mayor_precio = |$-\infty$|
    for i=1 to n
        precio = p[i] + corte_de_vara_recursivo(n-i, p)
        if precio > mayor_precio
            mayor_precio = precio
          
    return mayor_precio
\end{pseudocode}
El raciocinio que usamos para esta implementación deja claro que el problema cuenta con subestructura óptima. Falta ver si también presenta subproblemas superpuestos, pues solo en ese caso un algoritmo de programación dinámica mejoraría la eficiencia.

% TODO Por qué presenta problemas superpuestos? Árbol de recursión? O: explicar el ciclo, quizás eso es más fácil

\subsubsection{Solución top-down}

Para eliminar la redundancia que introducen los problemas superpuestos, se implementa un algoritmo top-down. Se añade una función principal que crea un arreglo donde la función auxiliar recursiva memoiza el cálculo del mejor precio para cada longitud.
\begin{pseudocode}
función corte_de_vara_top_down(n: entero, p: arreglo de enteros) entero
    |\br{arri}|precios_optimos = |\text{arreglo de enteros de longitud $n+1$}|
    precios_optimos[0] = 0

    for i=1 to n
        precios_optimos = |$-\infty$||\br{arrf}|

    return corte_de_vara_memoizado(n, p, precios_optimos)

función corte_de_vara_memoizado(
    n: entero, 
    p: arreglo de enteros, 
    precios_optimos: arreglo de enteros
) entero
    |\br{memoi}|if precios_optimos[n] >= 0
        return precios_optimos[n]|\br{memof}|

    mayor_precio = |$-\infty$|
    for i=1 to n
        precio = p[i] + corte_de_vara_memoizado(n-i, p, precios_optimos)
        if precio > mayor_precio
            mayor_precio = precio

    |\br{reti}|precios_optimos[n] = mayor_precio
    return mayor_precio|\br{retf}|
\end{pseudocode}
\begin{annotations}
  \codebrace[xshift=5.5cm, width=4.5cm]{arri}{arrf}{Se crea el arreglo lleno de $-\infty$ (salvo el caso base) para indicar las soluciones que aún no se han memoizado.}
  \codebrace[width=8cm]{memoi}{memof}{Si el valor está memoizado, se retorna enseguida para no recomputarlo.}
  \codebrace[xshift=3cm, width=5.7cm]{reti}{retf}{El valor se calcula normalmente y justo antes de retornarlo, se memoiza para los llamados siguientes.}
\end{annotations}

\subsubsection{Solución bottom-up}

Para diseñar el algoritmo bottom-up es necesario cambiar levemente de perspectiva. Tanto la solución recursiva como el algoritmo top-down ya contienen un ciclo, por lo que esta solución necesita dos ciclos: uno externo, que toma el lugar de la recursión y uno interno que sigue siendo muy similar al ciclo de los algoritmos previos.

\begin{pseudocode}
función corte_de_vara_bottom_up(n: entero, p: arreglo de enteros) entero
    |\br{arri}|precios_optimos = |\text{arreglo de enteros de longitud $n+1$}|
    precios_optimos[0] = 0|\br{arrf}|

    |\bxl{cic}|for longitud=1 to n|\bxr{cic}|
        |\br{inti}|mayor_precio = |$-\infty$|
        for i=1 to longitud
            precio = p[i] + precios_optimos[longitud-i]
            if precio > mayor_precio
                mayor_precio = precio

        precios_optimos[longitud] = mayor_precio|\br{intf}|
    
    return precios_optimos[n]
\end{pseudocode}
\begin{annotations}
  \codebrace[xshift=5.5cm, yshift=-0.2cm, width=5cm]{arri}{arrf}{En bottom-up, el arreglo se llena en orden: no hay que llenarlo de antemano como en top-down, con el primer elemento basta.}
  \codebrace[xshift=1.5cm, width=6cm]{inti}{intf}{El ciclo interno es prácticamente idéntico al de los algoritmos anteriores. Cambia que \texttt{n} ahora es \texttt{longitud} y el uso de la tabulación en lugar de un llamado recursivo}
  \codebox[xshift=1cm, width=3.8cm]{cic}{Este ciclo hace que se solucione el problema de corte de vara para cada valor de \texttt{longitud} de \texttt{1} a \texttt{n}}
\end{annotations}

Esa solución cuenta con una complejidad temporal de $\Theta(n^2)$, por los dos ciclos anidados, y con una complejidad espacial de $\Theta(n)$, por el arreglo que se utiliza para la tabulación.

\subsubsection{Obtener los tamaños de los cortes}

Inicialmente, el problema de corte de vara se planteó como un problema de optimización, donde la respuesta esperada era el máximo ingreso posible. Sin embargo, el problema se podría plantear como un problema de búsqueda en el que se busca alguna configuración de cortes que maximiza el ingreso, en lugar del ingreso en sí.
\begin{itemize}
  \item Entradas: Un entero $n$ que representa la longitud de la vara y un arreglo $p$ de tamaño $n+1$, donde $p[i]$ es el precio de una vara de longitud $i$.
  \item Salida: Una pila con las longitudes de las piezas que maximizan el precio de venta.
\end{itemize}

Para esto, al implementar el algoritmo de programación dinámica se debe no solo computar el valor óptimo sino también almacenar información para que, una vez se obtenga el valor óptimo, se pueda recuperar los pasos que se siguieron para obtenerlo. 

En este caso, además de almacenar el precio óptimo para una longitud $\ell$, se almacena el primer corte que corresponde a ese precio óptimo. Solo se almacena el primer corte $c$ para cada longitud, porque para determinar posteriormente el resto de cortes (si se requieren más de un corte), basta con examinar el subproblema de la longitud restante, o sea, el subproblema de la vara de longitud $\ell-c$ (para el cual también se guardó el primer corte).

\begin{pseudocode}
función corte_de_vara_backtracking(n: entero, p: arreglo de enteros) pila de enteros
    precios_optimos = \text{arreglo de enteros de longitud $n+1$}
    cortes_optimos = \text{arreglo de enteros de longitud $n+1$}
    precios_optimos[0] = 0

    for longitud=1 to n
        mayor_precio = |$-\infty$|
        for i=1 to longitud
            precio = p[i] + precios_optimos[longitud-i]
            if precio > mayor_precio
                mayor_precio = precio
                |\bxl{mejor}|mejor_corte = i|\bxr{mejor}|
        precios_optimos[longitud] = mayor_precio
        |\bxl{optm}|cortes_optimos[longitud] = mejor_corte|\bxr{optm}|

    segmentos = |\text{pila de enteros de longitud máxima $n$}|
    |\br{cortsi}|i = n
    while i > 0:
        corte = cortes_optimos[i]
        segmentos.put(corte)
        i -= corte|\br{cortsf}|

    return segmentos
\end{pseudocode}
\begin{annotations}
    \codebox[notepos={(cbox@mejor.east)+(5cm,0.6cm)}, curve={bend right=5}, width=4.7cm]{mejor}{Se almacena el primer corte que se hizo en la solución óptima de ese subproblema.}
    \codebox[xshift=2.7cm, width=5.7cm]{optm}{Para una longitud $\ell$, el primer corte es \inline{cortes_optimos[$\ell$]}; los cortes restantes (si hay) se deben averiguar mirando el subproblema de la longitud restante, que es \inline{$\ell - $cortes_optimos[$\ell$]}}
    \codebrace[xshift=3cm, width=6cm]{cortsi}{cortsf}{Se toma el primer corte óptimo para \inline{n}; se resta de \inline{n} para obtener una nueva longitud, luego se mira el primer corte para ella; y así hasta acabar la vara}
\end{annotations}

\practicar[gfg]{Rod Cutting}{https://www.geeksforgeeks.org/problems/rod-cutting0840/1}










